\section{Results}

\subsection{Evidence base and study selection}

The current workbook contained 64 coded review records representing 62 unique normalized titles. Publication years were 2025 (28); 2024 (21); 2023 (9); 2022 (6). The most frequently matched domains were Cardiovascular (32); Oncology (25); Surgery / orthopedics / dentistry (24); Neurology / neuroscience (18); Diabetes / endocrine (17); the leading function categories were Treatment / therapy planning (56); Prediction / prognosis / risk (49); Diagnosis / detection (46); Education / training (19); Personalized care (9). The most frequently matched barrier groups were Privacy / security / governance (51); Ethical / legal / regulatory (51); Computing / scalability / resources (51); Interoperability / standards (49). Counts are record-level, rule-based, and non-mutually exclusive.


The evidence-base characteristics are summarized in Table~\ref{tab:characteristics}. Figure~\ref{fig:prisma} reports the project selection counts. Automated arithmetic checking identified a five-record discrepancy between the reported full-text exclusions and included studies; the source values are retained pending reconciliation.

\begin{figure}[H]
    \centering
    \includegraphics[width=\columnwidth]{fig1_prisma_single_portrait.pdf}
    \caption{PRISMA 2020 flow of records and reports. Source counts are retained pending reconciliation of the full-text arithmetic.}
    \label{fig:prisma}
\end{figure}

\begin{table*}[htbp]
\centering
\small
\caption{Characteristics of the included review evidence base.}
\label{tab:characteristics}
\begin{tabularx}{\textwidth}{@{}>{\raggedright\arraybackslash}p{0.20\textwidth}YY@{}}
\toprule
\textbf{Dimension} & \textbf{Current extracted signal} & \textbf{Interpretive use} \\
\midrule
Evidence base & 64 coded rows; 62 unique normalized titles; 2 duplicate-title rows. & Defines the analytic set and makes title-level duplicate reconciliation explicit. \\
Publication years & 2022: 6; 2023: 9; 2024: 21; 2025: 28; missing: 0. & Shows the temporal distribution using the canonical publication-year field with documented fallback. \\
Protocol scope & English-language reviews from 2020 onward addressing patient-linked physiological, biomechanical, biological, or adjacent computational models. & Keeps eligibility interpretation tied to the PROSPERO protocol. \\
Dominant domains & Cardiovascular (32); Oncology (25); Surgery / orthopedics / dentistry (24); Neurology / neuroscience (18); Diabetes / endocrine (17); General / multisystem (15) & Provides the empirical basis for the application landscape. \\
Quality context & Quality-appraisal and reproducibility fields are heterogeneous and incompletely populated. & Supports descriptive synthesis and cautious interpretation rather than pooled certainty claims. \\
\bottomrule
\end{tabularx}
\end{table*}


\subsection{Conceptual definitions and the digital patient spectrum}

The included reviews used digital twin terminology alongside patient digital twin, human digital twin, virtual patient, in silico patient or trial, digital patient, synthetic patient, and adjacent constructs. Patient specificity, biological hierarchy, temporal updating, and coupling direction varied independently, supporting a multidimensional spectrum rather than a single lexical definition (Table~\ref{tab:concepts}; Figure~\ref{fig:spectrum}).

\begin{table*}[htbp]
\centering
\small
\caption{Conceptual definitions and terminology across the digital patient spectrum.}
\label{tab:concepts}
\begin{tabularx}{\textwidth}{@{}>{\raggedright\arraybackslash}p{0.20\textwidth}YY@{}}
\toprule
\textbf{Conceptual dimension} & \textbf{Current extracted signal} & \textbf{Interpretive use} \\
\midrule
Primary terminology & Digital twin (47); Virtual patient (5); Human digital twin (5); Patient digital twin (4); In silico patient / trial (3) & Shows the distribution of the first identifiable primary construct per review record. \\
Patient specificity & Patient-specific/individualized (56); generic/population/archetype (28). & Separates individualized models from population-level or archetypal constructs; categories may overlap. \\
Temporal orientation & Dynamic/continuous (54); explicitly static (15). & Tests whether digital twin terminology is accompanied by longitudinal updating. \\
Data coupling & Bidirectional (33); one-way (39). & Distinguishes model-, shadow-, and twin-like coupling claims. \\
Hierarchy & Hierarchy levels are converted to an ordinal cell-to-system-of-systems scale for Figure 2; multi-level entries use the highest coded level. & Makes the spectrum plot reproducible while preserving the row-level source text. \\
\bottomrule
\end{tabularx}
\end{table*}


\begin{figure*}[!tbp]
    \centering
    \includegraphics[width=0.96\textwidth]{fig2_spectrum_double_landscape.pdf}
    \caption{Digital patient spectrum. Records sharing a hierarchy--coupling category are deterministically spread within that cell to reduce overlap. Color denotes healthcare domain, marker shape denotes the primary construct, and marker area denotes the reported number of included papers.}
    \label{fig:spectrum}
\end{figure*}

\subsection{Modeling and computational pipeline}

Construction approaches were multimodal and heterogeneous. The extraction framework distinguishes data inputs, model families, personalization and coupling, named interoperability standards, computing infrastructure, and reported clinical or research outputs (Table~\ref{tab:pipeline}; Figure~\ref{fig:pipeline}).

\input{generated/tables/table3_pipeline}

\begin{figure*}[!tbp]
    \centering
    \includegraphics[width=0.96\textwidth]{fig3_pipeline_double_landscape.pdf}
    \caption{Computational pipeline across included reviews. Values are record-level mentions and categories are not mutually exclusive.}
    \label{fig:pipeline}
\end{figure*}

\subsection{Validation, fidelity, and clinical maturity}

Validation reporting ranged from conceptual discussion to empirical patient-data evaluation, simulation or benchmark comparison, cross-validation, and limited prospective or operational evidence. The narrative extraction also captured fidelity, uncertainty, reproducibility, quality appraisal, and maturity (Table~\ref{tab:validation}; Figure~\ref{fig:validation}).

\begin{table*}[htbp]
\centering
\small
\caption{Validation strategies and evidence strength.}
\label{tab:validation}
\begin{tabularx}{\textwidth}{@{}>{\raggedright\arraybackslash}p{0.20\textwidth}YY@{}}
\toprule
\textbf{Validation dimension} & \textbf{Current extracted signal} & \textbf{Interpretation} \\
\midrule
Validation approaches & Empirical / patient-data validation (25); Simulation / benchmark comparison (18); Sensitivity / uncertainty analysis (13); Cross-validation / holdout (8); Conceptual / no model evaluation (6); Expert evaluation (3) & Reports record-level mentions of empirical, simulated, cross-validation, uncertainty, expert, and conceptual approaches. \\
Validation data source & Patient/clinical, simulation/benchmark, and other sources remain free-text and are preserved in the record-level output. & Avoids imposing an unsupported pooled effect interpretation. \\
Fidelity and uncertainty & Sensitivity/uncertainty language appears in 13 records. & Shows how often model credibility was discussed explicitly. \\
Reproducibility & Explicit affirmative, code, open-source, or reproducibility language appears in 26 records. & Identifies an actionable transparency gap. \\
Clinical readiness & Prototype / proof of concept (43); Theoretical / conceptual (21); Deployed / operational (19); Prospective / clinical study (3); Retrospective validation (1) & Separates conceptual promise, prototypes, validation studies, and deployed systems. \\
\bottomrule
\end{tabularx}
\end{table*}


\begin{figure*}[!tbp]
    \centering
    \includegraphics[width=0.94\textwidth]{fig4_validation_double_landscape.pdf}
    \caption{Validation approaches and reported maturity. Bars show non-mutually exclusive record-level mentions.}
    \label{fig:validation}
\end{figure*}

\subsection{Applications}

Applications spanned multiple clinical and research domains. Because reviews commonly addressed more than one domain or function, the application landscape is represented as a co-occurrence matrix rather than a forced single-label classification (Table~\ref{tab:applications}; Figure~\ref{fig:applications}).

\input{generated/tables/table5_applications}

\begin{figure*}[!tbp]
    \centering
    \includegraphics[width=0.94\textwidth]{fig5_applications_double_landscape.pdf}
    \caption{Application landscape. Cells show the number of review records in which a healthcare domain and function co-occurred.}
    \label{fig:applications}
\end{figure*}

\subsection{Barriers and future directions}

Cross-cutting limitations were classified into interoperability and standards; privacy, security, and governance; ethical, legal, and regulatory issues; computing and scalability; clinical adoption and workflow; data quality and availability; and validation, evidence, and reproducibility (Table~\ref{tab:barriers}; Figure~\ref{fig:barriers}). The corresponding reporting and implementation priorities are summarized in Table~\ref{tab:recommendations}.

\begin{figure}[H]
    \centering
    \includegraphics[width=\columnwidth]{fig6_barriers_single_portrait.pdf}
    \caption{Cross-cutting barriers. Bars show non-mutually exclusive record-level mentions.}
    \label{fig:barriers}
\end{figure}

\input{generated/tables/table6_barriers}
\begin{table*}[htbp]
\centering
\small
\caption{Future directions and recommendations derived from the extraction framework.}
\label{tab:recommendations}
\begin{tabularx}{\textwidth}{@{}>{\raggedright\arraybackslash}p{0.20\textwidth}YY@{}}
\toprule
\textbf{Recommendation} & \textbf{Priority} & \textbf{Evidence basis and suggested action} \\
\midrule
Standardize terminology & High & Report explicit criteria for digital model, digital shadow, digital twin, patient digital twin, virtual patient, and in silico patient. \\
Define maturity levels & High & Classify patient specificity, hierarchy, temporal updating, data coupling, feedback, validation, and clinical integration. \\
Strengthen validation & High & Require clear validation datasets, internal/external validation where relevant, and sensitivity or uncertainty analysis. \\
Build interoperable infrastructure & High & Report applicable FHIR, HL7, DICOM, OMOP, ontology, identifier, and governance mechanisms. \\
Provide reproducible supplements & High & Publish row-level extraction outputs, executable rules, data-quality checks, and plotting datasets. \\
Study workflow and governance & Medium & Evaluate clinician trust, liability, consent, bias, interpretability, equity, and implementation burden. \\
Move toward prospective evidence & High & Distinguish conceptual promise from retrospective validation, prospective studies, and deployed clinical workflows. \\
\bottomrule
\end{tabularx}
\end{table*}


\FloatBarrier
